\apendices

\chapter{Índice lógico dos pipelines de análise}
\label{ap:pipelines}

\notaguia{Fonte: \texttt{Textos/indice\_logico\_pipelines.md}. Tabela
perna/fase/papel dos 54 pipelines --- o número do pipeline é identidade
estável e não deve ser renumerado.}

\chapter{Registro de decisões metodológicas (D1--D27)}
\label{ap:decisoes}

\notaguia{Fonte: seção ``As vinte e sete decisões'' da visualização. Cada
decisão: contexto, alternativas, escolha e consequência.}

\section{Decisões metodológicas detalhadas}
\label{ap:decisoes-detalhe}

\subsection{Censo de pixels em lugar de amostra (D21--D24)}
\label{ap:d21}

A análise da idade da pastagem no momento da conversão baseava-se
originalmente numa amostra de pixels sorteada sobre o retângulo envolvente
do estado --- e não sobre o polígono estadual. Como Goiás tem forma
irregular, o retângulo capturava faixas de estados vizinhos: a verificação
por ponto em polígono mostrou que 43,7\% das linhas coletadas caíam fora do
território goiano. O filtro corretivo é estatisticamente válido --- condicionada
a cair em Goiás, a seleção permanece uniforme --- mas custava tamanho de
amostra e deixava intacta a causa: a fração efetivamente goiana variava de
32,8\% a 79,1\% conforme o ano, e essa fração \emph{caía} ao longo da série,
o que compromete qualquer estatística agregada por período.

A validação do procedimento revelou um segundo defeito, mais grave. A classe
21 do MapBiomas --- ``Mosaico de Usos'' --- estava ausente do dicionário de
grupos, e o código atribuía a categoria de ``censurado'' a qualquer classe
não encontrada. Uma idade perfeitamente conhecida era, assim, descartada
como desconhecida. A taxa de censura publicada, de 74,9\%, era na verdade de
63,7\%, e o número de observações utilizáveis passou de 11.035 para 15.933,
um aumento de 44\%. O ponto que torna esse defeito sério não é a magnitude:
é que os registros descartados não eram aleatórios. Eram justamente os de
origem mista entre agricultura e pastagem --- a categoria substantivamente
mais próxima de um dos dois mecanismos que a análise se propunha a
distinguir. O erro removia evidência preferencialmente de um lado da
conclusão. A causa-raiz é de projeto: o código usava o mesmo valor para
``censurado'' e para ``busca no dicionário falhou'', que são coisas
diferentes e exigem códigos diferentes (decisão D22).

A reconstrução substituiu a amostra por um censo: o cubo completo de
transições do estado foi processado localmente, resultando em 44.639.028
eventos de conversão --- mil vezes o tamanho da amostra --- cobrindo
3.817.080 ha, ou 11,2\% do território, em 244 dos 246 municípios. A
comparação entre amostra e censo foi desenhada em dois níveis, porque as
duas fontes divergem por causas independentes que o agregado confunde: ano a
ano, a amostragem mostrou-se sadia (diferença média de mediana de $-0{,}09$
ano); no agregado por período, as diferenças eram inteiramente de
composição, isto é, de peso relativo dos anos dentro do período. Daí a
decisão D24, que estabelece o contrato de ponderação: um módulo que aceita
dado censitário deve ser explicitamente ciente de peso, sob pena de o
gradiente estimado desaparecer sem aviso.

\subsection{A mudança de rótulo do Mosaico de Usos (D25--D26)}
\label{ap:d25}

No trecho final da série --- de modo agudo entre 2022 e 2024 --- a conversão
de pastagem em agricultura passa a ser crescentemente rotulada como classe
21, ``Mosaico de Usos''. A consequência é que qualquer medida dependente da
classe ``Agricultura'' subconta a lavoura recente numa janela que toca
2020--2024. Não se trata de interrupção do fenômeno: a estatística de área
plantada de soja registra crescimento de 38\% no mesmo período. Foi o rótulo
que mudou.

A tentação natural --- somar agricultura e mosaico --- não é uma correção,
e o trabalho recusa tratá-la como tal. Somar as duas classes equivale a
supor que todo o Mosaico é agricultura mal rotulada, o que é falso: a classe
contém também sistemas integrados de lavoura e pecuária, mosaicos finos
demais para separar a 30~m e mosaico antigo e estável, existente desde muito
antes do fenômeno. A união, portanto, superconta tanto quanto a classe
isolada subconta.

A decisão D26 estabelece o enquadramento correto: em vez de um ponto,
reporta-se o intervalo

\begin{equation*}
\big[\,\text{agricultura}\,;\ \text{agricultura} \cup \text{mosaico}\,\big],
\end{equation*}

\noindent cujo extremo inferior subconta e cujo extremo superior superconta,
considerando-se robusta apenas a conclusão que sobrevive nos dois extremos. Se um resultado muda de sinal ou de significância entre os
extremos, ele depende da convenção de classe --- e isso precisa ser dito.

Duas qualificações completam a decisão. Primeiro, a régua superior responde
a uma pergunta legítima, embora diferente e mais grossa: não ``quanta terra
virou lavoura?'', mas ``quanta terra saiu de pastagem para lavoura ou uso
misto?'' --- pergunta robusta à reetiquetagem por construção, já que mover
um pixel entre as duas classes não altera a soma. O erro seria passar uma
pergunta pela outra. Segundo, o intervalo é ferramenta de contenção do
artefato, não a melhor evidência disponível: onde existe fonte independente
do classificador --- área plantada, rebanho, crédito, valor adicionado ---,
é ela que responde, e o intervalo fica reservado ao que só o sensoriamento
remoto mede.

\subsection{Oscilação classificatória entre pastagem e savana}
\label{ap:oscilacao}

Nas matrizes de transição, o fluxo reverso de pastagem para vegetação
natural é volumoso, e a leitura tentadora seria de regeneração. A
verificação em classe bruta, sem colapsar os grupos, mostra que esse fluxo
é dirigido quase inteiramente à formação savânica --- de 75\% a 98\%
conforme o par de anos --- e praticamente nunca à formação campestre,
contrariando a hipótese inicial de que a confusão se daria entre pastagem e
campo natural.

O diagnóstico decisivo é a simetria: a razão entre o fluxo savana para
pastagem e o fluxo inverso colapsa de 8,4 vezes, no primeiro par decenal,
para 1,22 vezes no par anual mais recente. Num único ano registram-se
dezenas de milhares de hectares em cada direção, quase equilibrados.
Regeneração real é lenta e unidirecional; oscilação de classificador é
bidirecional e balanceada. Há um componente real minoritário, em formação
florestal, que só se acumula em horizontes longos. A consequência
operacional é uma proibição de leitura, adotada em todo o texto: o fluxo
reverso de pastagem para vegetação natural nunca é narrado como recuperação.
As conclusões centrais do trabalho não dependem dele --- apoiam-se no fluxo
direto de conversão e em saldos líquidos, ambos robustos a uma oscilação
simétrica, que se cancela no líquido.

\section{Resultados superados}
\label{ap:superados}

O trabalho registra quatro resultados que teriam sido publicados
se um teste específico não tivesse sido feito. Nenhum deles era ingênuo: em
cada caso, a versão superada era a leitura mais plausível com o que havia
até ali. Em dois casos ela afirmava mais do que a versão que ficou no lugar;
nos outros dois, contrariava a tese. Os resultados corrigidos são os
apresentados no Capítulo~\ref{cap:resultados}; o que a
Quadro~\ref{quadro:superados} registra é a distância entre as duas versões.

\begin{quadro}[htb]
\centering
\caption[Resultados superados e o teste que os superou]{Resultados superados e o teste que os superou.}
\label{quadro:superados}
\small
\begin{tabular}{p{4.2cm}p{4.6cm}p{4.4cm}}
\toprule
\textbf{Versão superada} & \textbf{O que o teste mostrou} & \textbf{Decisão} \\
\midrule
A agricultura desacelerou depois de 2019 &
Mudou o rótulo, não o território: as medidas de campo marcham na mesma
janela &
D25--D26 \\
\addlinespace
O Norte antecede o Sul no tempo &
Regressão espúria; sob método robusto à integração, a precedência some nas
duas direções &
D16 \\
\addlinespace
O \emph{drive} comum é estatisticamente significante &
Inferência apropriada ao desenho leva o p de $\sim$0,03 a 0,07--0,13 &
D20 \\
\addlinespace
A marcha ao norte sem barra de erro &
Com IC por \emph{bootstrap}, a vegetação natural passa a ``ancorada'' &
D19 \\
\bottomrule
\end{tabular}
\fontefig{elaboração própria.}
\end{quadro}

A esses somam-se, apresentados nas frentes correspondentes, o gradiente
latitudinal de idade da pastagem, o plantio direto como mecanismo e a
leitura da vegetação como bloco homogêneo ao norte. A auditoria de cada um
está registrada junto ao resultado que ela corrigiu.

\chapter{Glossário de métricas}
\label{ap:glossario}

\notaguia{Fonte: \texttt{Textos/metodologia/glossario\_metricas.md} + glossário
da visualização.}
